\section{Task suite and protocol}
\label{sec:tasks}

The suite has \NTasks{} tasks: \NTaskAlgorithms{} algorithms, each with two variants
(Appendix~\ref{app:tasks}). A task states a rule in prose with its parameters. The model
must implement the rule as \texttt{solve(prices, split\_factors)}, returning a date-by-asset
frame. For portfolio and signal tasks the frame holds positions. For volatility, risk and
forecast tasks it holds the value known at the previous close. Each task can be solved
either by hand or by one library call (\texttt{positions()} or \texttt{rolling()}).

\paragraph{Data.}
The inputs are synthetic four-asset panels of \NRows{} business days with seeded Gaussian
returns. Raw quotes are adjusted closes divided by cumulative split factors. A panel has four
splits at fixed rows: 2-for-1, 3-for-1, 3-for-2 and a 1-for-2 reverse split. The
\NEFourTasks{} tasks carried over from our earlier study keep \NRowsEFour{} rows, and with them
only the first two splits, so they reproduce that study's inputs exactly. Variant~a gives cumulative factors in alphabetical column order.
Variant~b gives event ratios, a shuffled column order and different parameters. Every task
runs on \NSeeds{} data seeds, and held-out tasks also run on one real
panel (below). Together, the tasks contain the failures we observed in earlier
agent-written programs:
\begin{itemize}
\item cumulative factors applied again, or event ratios used as cumulative factors;
\item positions decided on the close they trade over;
\item warm-up rows computed on a partial window or left missing;
\item the missing first return of \texttt{pct\_change()};
\item columns assumed to be in alphabetical order.
\end{itemize}

\paragraph{Real prices.}
The real panel holds the as-traded daily closes of Amazon, Shopify, Alphabet (class~A) and
Tesla for the \NRealRows{} sessions from \RealFirst{} to \RealLast{}. All four split in that
window: 20-for-1, 10-for-1, 20-for-1 and 3-for-1 (Appendix Figure~\ref{fig:real}).
Amazon's split falls on row
\NRealSplitRow{}, where asset~A splits in the synthetic panels. Columns carry the tickers, in
split order for variant~a and shuffled for variant~b, and each variant keeps its factor
convention. The closes come from Yahoo
Finance, which serves them split-adjusted. We rebuild the as-traded quotes from a split table
checked against the companies' announcements, and the rebuilt closes before each split match
published figures. Yahoo's terms forbid redistribution, so we release the fetch code and the
file's SHA-256 rather than the prices. The panel has no reverse split and no dividend.

\paragraph{References.}
Each task has an independent reference written with numpy, which does not call \lib{}. On
every task and seed, the reference agrees with the one-step library call to within
$\RefGapOther$. The exception is \texttt{min\_variance}, which agrees to within
$\RefGapMinVariance$: the library solves it numerically, and the reference enumerates the
supports of the long-only solution exactly. Building these references exposed a precision
defect in the library's long-only minimum-variance optimiser. On daily-scale covariances the
optimiser stopped early, and weights differed from the exact solution by up to
$1.2\times10^{-2}$. Rescaling the covariance before solving fixed it.

\paragraph{Judge and hidden grader.}
\label{sec:grader}
The judge runs each program in a separate process with a memory limit and a 120-second
timeout. In our runs the process is also confined on Linux with Landlock and seccomp, which
restrict its file-system and network access. Without library access, the judge refuses
\texttt{import fin\_skills}. It applies the three checks of \S\ref{sec:checks}, perturbing
from rows 32 and 64 for \texttt{execution\_lag} (Figure~\ref{fig:episode}c). The grader is
separate code that never
appears in a prompt. A program is \emph{valid} only if it ran, passed every check, and
matches the reference within an absolute tolerance of $10^{-8}$ ($10^{-5}$ for
\texttt{min\_variance}). A program that passes every check but is wrong is \emph{accepted but
wrong}; we report this rate because it is the error the checks cannot see.

Before any model was run, the suite had to pass three qualification tests:
\begin{itemize}
\item on every task, a correct program (the one-step library call) passes every check and is
graded valid, so no check raises a false alarm;
\item on every task and seed, the reference agrees with the library;
\item on a position task and a table task, a planted version of each failure fails both its
check and the grade: trading on the same close, no adjustment, factors cumulated twice
(position task only), and missing warm-up rows.
\end{itemize}
All three also pass on the real panel; the first is run on the held-out tasks, the only ones
the real panel is used for.

\paragraph{Development and held-out split.}
The split is by algorithm, so no held-out algorithm appears during development. The
\NForcedDev{} algorithms that a model has already seen go to development: four from our
earlier study and one used in model qualification. Of the remaining \NDrawn{}, a seeded draw
sends \NDrawnDev{} to development. Development then has \NDevAlgorithms{} algorithms
(\NDevTasks{} tasks). The held-out set has \NHeldAlgorithms{} algorithms (\NHeldTasks{} tasks)
and covers all \NFamilies{} families.

\paragraph{Protocol.}
The primary outcome is the valid-delivery rate on held-out tasks. Secondary outcomes are
the accepted-but-wrong rate, model calls, generated tokens, wall time and library use.

The main comparison runs \NMethods{} methods $\times$ 2 arms (\texttt{cards},
\texttt{cards+library}) $\times$ \NHeldTasks{} tasks $\times$ \NSeeds{} seeds $\times$
\NModels{} models, or \NMainEpisodes{} episodes. The same comparison on the real panel adds
\NRealEpisodes{} episodes, which are reported separately and never pooled with the synthetic
seeds. The component ablation adds
\texttt{NoImprove}, \texttt{CheckFeedback} and \texttt{PlaybookLearning} on the
\texttt{none} and \texttt{library} arms, a further \NAblationEpisodes{} episodes.

Intervals are task-clustered bootstrap intervals, since seeds of one task are not
independent. Paired comparisons match the task, seed and model, with a Holm correction. The
analysis plan is frozen in a protocol file before inference, and every result is reported
whatever its direction. \pending{Models: \NModels{} open-weight models from at least four
families, chosen by qualification on development tasks; to be named when the protocol is
frozen.}
